\begin{definition}[A nonabelian dual seam word]
    \label{def:peps_dual_nonabelian_seam_word}
    \lean{TNLean.PEPS.torusDualSeamWord}\leanok
    Choose elements $u,v$ in any coefficient group. Traverse a labelled dual path
    in order, multiplying by $u$ on eastward horizontal seam crossings,
    by $v$ on northward vertical seam crossings, and by their inverses
    on the reverse crossings. This is a homotopy detector, independent
    of the physical flux insertion assigned to the path. If $E_{x,y},N_{x,y}$
    are positive steps, define
    \begin{align}
        \omega(E_{x,y})=u^{[x+1=0]},\qquad
        \omega(N_{x,y})=v^{[y+1=0]},\qquad
        \omega(e^{-1})=\omega(e)^{-1}.
        \label{eq:peps_dual_seam_step_word}
    \end{align}
    The brackets denote the zero-or-one indicator. The path word satisfies
    $W(1_a)=1$ and $W(pe)=W(p)\omega(e)$.
\end{definition}

\begin{lemma}[Restricted moves preserve the nonabelian seam word]
    \label{lem:peps_dual_restricted_seam_word}
    \lean{TNLean.PEPS.TorusDualHomotopy.seamWord_eq_of_origin_not_mem}\leanok
    If the allowed swept vertex set excludes the origin, every generated
    homotopy preserves the seam word for every coefficient group and every
    choice of its two generators.
\end{lemma}
\begin{proof}\leanok
    \uses{def:peps_dual_nonabelian_seam_word,def:peps_dual_flux_homotopy}
    Concatenation multiplies the words, reversal inverts them, and
    backtracks cancel. Put $\epsilon=[x+1=0]$ and $\delta=[y+1=0]$.
    On a square with support $(x+1,y+1)\ne(0,0)$, at least one indicator
    vanishes, so
    \begin{align}
        W(\mathrm{EN}_{x,y})=u^\epsilon v^\delta
        =v^\delta u^\epsilon=W(\mathrm{NE}_{x,y}).
        \label{eq:peps_restricted_seam_square}
    \end{align}
    The only excluded case is the square surrounding the origin.
\end{proof}

\begin{lemma}[Zero displacement does not classify restricted homotopy]
    \label{lem:peps_dual_restricted_displacement_obstruction}
    \lean{TNLean.PEPS.torusDualSquareLoop_obstructs_restricted_completeness,
        TNLean.PEPS.torusDualSquareLoop_not_homotopic_avoiding_origin}\leanok
    If a coefficient group has noncommuting elements $u,v$, the dual square loop
    surrounding the origin has zero lifted displacement, but it cannot
    be contracted using square moves supported away from that origin.
    The unconditional version constructs two noncommuting transpositions
    on three elements internally. This holds without restrictions on the
    two torus periods or on a physical flux group.
\end{lemma}
\begin{proof}\leanok
    \uses{lem:peps_dual_restricted_seam_word,def:peps_dual_flux_displacement,
        def:peps_dual_nonabelian_seam_word}
    The four displacements sum to zero, while the seam word is
    $uvu^{-1}v^{-1}\ne1$. The constant path has word $1$, so the preceding
    invariant rules out the asserted contraction.
\end{proof}
